\documentclass[../main]{subfiles}
\begin{document}
\chapter{Costos Fijos}
\img{images/01/costo.jpg}{Costos}
\section{Costos Fijos}
\info{¿Qué son los costos fijos?}{
  Los costos fijos son aquellos cuyo pago es asumido por la empresa de
  manera constante, independientemente de su participación dentro del
  proceso productivo. A estos costos se les conoce como fijos porque no
  varían ante los cambios de la producción de bienes y servicios. Es
  decir, son aquellos que sin importar cuánto se produzca, siempre
  deberán ser abonados. Por ejemplo, el alquiler de una oficina o un
  local, los sueldos administrativos, los servicios de telefonía e
internet, el pago de seguros, etc.}
\subsection{Como determinarlos}
Los costos fijos pueden determinarse acorde a las decisiones de la
administración, de tal forma que esta puede prescindir de un servicio
o por el contrario adquirir uno, o preservar los ya existentes según
sean las necesidades de la organización.
\subsection{Importancia}
La importancia de llevar control de estos costos radica en su gran
incidencia en la determinación de los precios finales a los bienes o
servicios que pueda proporcionar una empresa. Además de que deben ser
incluidos dentro de los presupuestos por ser costos que se presentan
de forma constante y permanente
Características de los costos fijos
A continuación, se relacionan las principales características de los
costos fijos:
\subsection{Características}
\begin{enumerate}
  \item No varían por la presencia de cambios en la producción.
  \item Pueden determinarse en función de las decisiones administrativas.
  \item A pesar de estar ligado a la capacidad de producción que
    posee una organización, estos no sufren variaciones si se llegan
    a alterar los niveles de producción.
  \item Cumplen un rol fundamental para el cálculo del punto de equilibrio.
    Se debe llevar registro detallado de estos costos para poder
    tener control oportuno sobre ellos.
\end{enumerate}
\begin{center}
  \begin{tikzpicture}
    \draw (0,0)--(0,5)[-stealth] node [left] {C};
    \draw (0,0)--(5,0)[-stealth] node [below] {Q};
    \draw (0,2)--(5,2) node [below] {CF:Costos Fijos};
  \end{tikzpicture}
\end{center}
\subsection{Tipos de costos fijos}
\subsubsection{Discrecionales}
Aquellos que pueden llegar a reducirse sin que la producción de la
organización sea afectada. Estos costos se originan con la finalidad
de ser empleados en actividades o áreas determinadas y las
organizaciones suelen ser más flexibles sobre estos. Los costos
discrecionales generalmente se originan en planes tácticos de
corto alcance generalmente para un marco tiempo de un año o más.
Por ejemplo: campañas de publicidad, los alquileres de ciertos
bienes, investigaciones, entre otros.
\subsubsection{Comprometidos}
Aquellos costos que no son susceptibles de sufrir modificaciones
porque el proceso productivo de la empresa puede ser afectado. Por
ende, los costos fijos comprometidos corresponden a costos necesarios
y de carácter obligatorio ya que, si estos son interrumpidos, pueden
presentarse inconvenientes severos en la rentabilidad de la
organización. Por el contrario, la interrupción de los procesos de
producción por algún factor externo a estos gastos, no alterará los
costos fijos comprometidos, debido a la naturaleza que presentan.

Los costos comprometidos se originan en planificaciones estratégicas
que generalmente abarcan un marco temporal de varios años y que en
gran medida fijan el nivel y el tipo de operaciones de una empresa.
Estos costos se deben controlar de manera cuidadosa con un adecuado
análisis previo, para optimizar de la manera más eficiente los
elementos que integran los costos fijos comprometidos.

Por ejemplo: salarios de empleados, seguros, impuestos sobre raíz de
propiedad, entre otros.
\subsection{Ejemplos de costos fijos}
A continuación se relacionan algunos ejemplos de costos fijos:
\begin{itemize}
  \item  El pago del alquiler de un local, bodega u oficina, que se
    realiza de forma puntual y sucesiva acorde a las condiciones
    estipuladas en el contrato.
  \item Pagos por seguros en caso de siniestros.
    Depreciación de maquinaria y equipo.
  \item Costos de investigación o desarrollo.
  \item Pagos por servicios relacionados no medibles de telefonía,
    internet, agua, etc.
  \item Amortización de costos de patentes, marcas y licencias.
  \item Los permisos municipales y licencias también son costos fijos.
  \item Los salarios de los empleados administrativos, no
    relacionados en cantidad con la cantidad de producción.
\end{itemize}
\actividad{Actividad en proyecto: Identificación de costos}
\begin{enumerate}
  \item Identificar los costos fijos en su proyecto, elabore un
    listado de todos los costos fijos anuales y mensuales, y
    colóqueles en una tabla.
  \item Justificar en cada caso porqué es un costo fijo.
\end{enumerate}
\actividad{Investigar y contestar el siguiente cuestionario}
\begin{enumerate}
  \item ¿Qué es la depreciación y explicar con sus palabras?
  \item ¿Qué es la amortización y los tiempos para rodados,
    inmuebles, muebles y máquinas herramientas?
  \item ¿Qué es un \emph{bien de capital}?
  \item ¿Qué es investigación y desarrollo?
  \item ¿Qué es un siniestro y porque es necesario pagar un seguro,?
    (ej.:seguro de vida)
  \item ¿Qué es una patente, licencia o permiso y porqué debe pagarse?
  \item ¿Cuáles son las tareas o trabajos que realizan los empleados
    administrativos?
  \item ¿Qué es una campaña de publicidad?
  \item ¿Qué es un proceso productivo?
  \item Explique que es el punto de equilibrio.
  \item ¿Qué es un registro de costos?
\end{enumerate}
\end{document}
